<<<<<<< HEAD
%% Capítulo 2 — Estado da Arte
\chapter{Estado da Arte}
\label{chap:estado_arte}

\textbf{OIDC / OAuth~2.0} — O \gls{oauth2} \cite{RFC6749} é a \textit{framework} de autorização padrão da indústria. O \gls{oidc} \cite{OIDCCore} acrescenta uma camada de identidade, introduzindo o \textit{ID Token} e o \textit{UserInfo Endpoint}. O fluxo recomendado para aplicações \textit{web} é \textit{Authorization Code com PKCE} \cite{RFC7636}; em testes usa-se \gls{ropc}.

\textbf{JWT e JWKS} — O \gls{jwt} \cite{RFC7519} é um token compacto com três partes (header, payload, signature) codificadas em Base64URL. O algoritmo RS256 usa criptografia assimétrica: o \gls{idp} assina com a chave privada e os \textit{resource servers} validam localmente com a chave pública obtida via \gls{jwks} \cite{RFC7517}, sem chamadas adicionais ao \gls{idp}.

\textbf{RBAC} — O \gls{rbac} \cite{Ferraiolo1992} associa permissões a papéis; os utilizadores herdam permissões por atribuição de papel. Simplifica a administração face a listas de controlo de acesso (\gls{acl}) e é adequado para a maioria dos cenários empresariais.

\textbf{MFA / TOTP} — A \gls{mfa} \cite{NIST80063B} exige dois ou mais fatores independentes. O \gls{totp} \cite{RFC6238} gera códigos de 6 dígitos, válidos 30 segundos, baseados numa chave partilhada e no \textit{timestamp}. O NIST SP~800-63B recomenda AAL2 (com \gls{mfa}) para acesso a dados sensíveis.

\textbf{JML} — O ciclo de vida Joiner–Mover–Leaver \cite{IDPro} automatiza o provisionamento, a transição e a revogação de acessos, eliminando contas órfãs e garantindo o princípio do privilégio mínimo em cada fase.
=======
%% ============================================================
%%  Capítulo 2 — Estado da Arte
%% ============================================================
\chapter{Estado da Arte}
\label{chap:estado_arte}

\section{Gestão de Identidade e Acesso}

A \gls{iam} é o conjunto de políticas, processos e tecnologias que garante que os utilizadores certos têm acesso aos recursos certos, no momento certo e pelos motivos certos \cite{NIST80063B}. Os pilares de um sistema \gls{iam} são:

\begin{itemize}
  \item \textbf{Autenticação}: verificar a identidade declarada pelo utilizador;
  \item \textbf{Autorização}: determinar o que o utilizador pode fazer;
  \item \textbf{Gestão do ciclo de vida}: criar, atualizar e desativar identidades de forma controlada;
  \item \textbf{Auditoria}: registar eventos para rastreabilidade, conformidade e deteção de ameaças.
\end{itemize}

\section{OpenID Connect e OAuth 2.0}

O \gls{oauth2} \cite{RFC6749} é uma \textit{framework} de autorização que permite que aplicações obtenham acesso limitado a contas de utilizadores, sem expor as credenciais. O \gls{oidc} \cite{OIDCCore} acrescenta uma camada de identidade sobre \gls{oauth2}, introduzindo o conceito de \textit{ID Token} e o \textit{UserInfo Endpoint}, que permitem que clientes verifiquem a identidade do utilizador e obtenham informação de perfil básica.

Os fluxos \gls{oidc} mais relevantes são:
\begin{itemize}
  \item \textbf{Authorization Code com PKCE} \cite{RFC7636} — recomendado para aplicações \textit{web} e \textit{mobile}; o \gls{pkce} protege contra ataques de interceção do código de autorização;
  \item \textbf{Resource Owner Password Credentials} (\gls{ropc}) — utilizado apenas em testes e ambientes controlados; transmite credenciais diretamente ao servidor de autorização.
\end{itemize}

\section{JSON Web Token (JWT)}

O \gls{jwt} \cite{RFC7519} é um standard aberto que define um método compacto para transmitir informação entre partes como um objeto \gls{json} verificável. É composto por três partes codificadas em Base64URL: \textit{header} (algoritmo), \textit{payload} (\textit{claims}) e \textit{signature}.

O algoritmo \textbf{RS256} utiliza criptografia assimétrica: o \gls{idp} assina com a chave privada e os \textit{resource servers} validam com a chave pública. A distribuição segura das chaves públicas é feita via \gls{jwks} \cite{RFC7517} — um \textit{endpoint} \gls{json} que publica o conjunto de chaves ativas do \gls{idp}.

\section{Controlo de Acesso Baseado em Papéis (RBAC)}

O \gls{rbac} \cite{Ferraiolo1992} é um modelo de controlo de acesso em que as permissões são associadas a \textbf{papéis}, e os utilizadores adquirem permissões ao serem atribuídos a esses papéis. Este modelo simplifica a administração em comparação com listas de controlo de acesso (\gls{acl}) e é adequado para a maioria dos cenários empresariais.

O NIST define quatro modelos \gls{rbac} progressivos: \textit{Core RBAC}, \textit{Hierarchical RBAC}, \textit{Constrained RBAC} e \textit{Symmetric RBAC}. Neste projeto adota-se o \textit{Core RBAC} com elementos de hierarquia implícita (o papel \texttt{admin} acede a tudo o que o \texttt{colaborador} acede).

\section{Autenticação Multifator (MFA) e TOTP}

A \gls{mfa} \cite{NIST80063B} exige que o utilizador prove a sua identidade com dois ou mais fatores independentes: \textit{algo que sabe} (password), \textit{algo que tem} (dispositivo \gls{totp}) e/ou \textit{algo que é} (biometria).

O \gls{totp} \cite{RFC6238} gera códigos numéricos temporários baseados numa chave partilhada (\textit{secret}) e no \textit{timestamp} atual. A especificação define uma janela de 30 segundos e códigos de 6 dígitos, compatível com Google Authenticator, Microsoft Authenticator e Authy.

O NIST SP~800-63B classifica os autenticadores em três níveis (\textit{Authentication Assurance Levels} — AAL1, AAL2, AAL3) e recomenda AAL2 (com \gls{mfa}) para acesso a dados sensíveis.

\section{Ciclo de Vida de Identidades — JML}

O ciclo de vida \gls{jml} refere-se aos três processos chave de gestão de identidades empresariais \cite{IDPro}:

\begin{itemize}
  \item \textbf{Joiner} (\textit{Integração}): provisionar o acesso inicial quando um utilizador entra na organização ou muda de função — o princípio do \textit{least privilege} deve ser aplicado desde o início;
  \item \textbf{Mover} (\textit{Transição}): ajustar os acessos quando um utilizador muda de papel ou departamento — a revogação do acesso anterior deve ser imediata;
  \item \textbf{Leaver} (\textit{Saída}): desprovisionar todos os acessos quando um utilizador sai — contas órfãs são um vetor de ataque crítico.
\end{itemize}

A automação destes processos através de \gls{api}s de administração (como a Keycloak Admin \gls{rest} \gls{api}) elimina o erro humano e garante consistência.
>>>>>>> df2eb5220b31d8e704cf4f670cff764b26edabce
